\begin{statement}{E.40}{Proved}{All retained Ferrari points pass full periodic-braid conjugacy}
For every one of the $864$ actual joint points of E.24,
with multiplier residue $1$ or $73$ modulo $144$, put
\[
 d=abc,\qquad d'=a'b'c',\qquad z_4=d^4
                    \quad\text{in }T_{23}=T_{2,6}\times T_{3,6}.
 \tag{E.40.1}
\]
There is an element $h\in T_{23}$ with
\[
                           d'=hd^\lambda h^{-1}.
 \tag{E.40.2}
\]
All compatible higher lifts of the displayed multiplier
residue satisfy the same finite comparison.  Together
with E.35 this verifies both full periodic tests in
E.39 at every retained point.

In the binary factor, the $96$ branch and line inputs
give two distinct values of $d'$.  For each of them,
\[
       d'z_4^j\sim_{T_{2,6}}d
           \quad\Longleftrightarrow\quad j=0\text{ or }2
                  \quad(j\in\mathbb Z/4).
 \tag{E.40.3}
\]
Four literal braid conjugators of lengths from zero
to $11088$ certify the positive comparisons.
In the ternary factor the $216$ points give fifteen
distinct values of $d'$, with
\[
       d'z_4^j\sim_{T_{3,6}}d
           \quad\Longleftrightarrow\quad j=0
                  \quad(j\in\mathbb Z/9).
 \tag{E.40.4}
\]
The fifteen saved conjugators have lengths from
six to $131$.  Every input and every positive
identity has been checked by a separate literal
word evaluator importing no conjugacy solver.
\end{statement}
\begin{proof}
We keep the complete finite target models and
the actual source reconstruction maps of E.24.
The new binary element has a nonidentity image
in the actual quotient $T_{2,6}\to S_3$.
Accordingly the conjugacy algorithm uses, for
each of the six recorded coset representatives $g$,
the necessary base equality
\[
                 \overline{d'}=\overline{g d g^{-1}}.
 \tag{E.40.5}
\]
Only the cosets satisfying this equality proceed
to the kernel calculation.  Above each such coset,
the residual difference lies in the actual
forty-dimensional collected two-group $U$.

The weight-by-weight centralizer method of E.32
then applies with $d$ in place of $d^2$.
At weight $r$, every element of the current
centralizer centralizes the comparison modulo
weight $r$.  Its commutator with that comparison
in the next elementary abelian layer is an
actual homomorphism: conjugation by $U$ acts
trivially on that layer, so the product
commutator identity becomes addition there.
All the kernel basis vectors are retained and
closed under the actual two-group multiplication.
The computed subgroup orders verify the index
of the next centralizer as exactly the image
order of this homomorphism.  Thus every possible
conjugator above every matching base coset is
retained.  An inconsistent system eliminates
the whole corresponding coset, and all six
base cosets are checked.

As an exhaustive independent control, use the
degree-two target, of order $1536$, with the
same nonidentity-base element $abc$.
Enumerate all elements and all of its conjugates.
The conjugacy class has $96$ elements, and all
$1536$ decisions of the new algorithm agree
with this complete enumeration.

At degree six, reconstruct each of the $24$
original binary inputs, both multiplier branches,
and both parameters in its full correction line.
This gives all $96$ inputs in the statement.
Apply the complete algorithm to all four central
translates of each distinct value.  The outcome
is exactly \textup{(E.40.3)}.
The forty collected kernel-basis elements are
rebuilt from the saved original braid-word
expressions, including every square and commutator.
Reducing each successful conjugator in this
basis and adjoining its actual outer word
produces the four literal witnesses in the
statement.  Substitution verifies each identity
in the original finite group.

In the binary target, $|d|=16$ and $|z_4|=4$.
For the actual multiplier residues,
\[
 d^\lambda=d z_4^{(\lambda-1)/4},\qquad
 -\frac{\lambda-1}{4}\equiv
       \begin{cases}0&\lambda=1,\\2&\lambda=73\end{cases}
                                          \pmod4.
 \tag{E.40.6}
\]
Both required central exponents occur in
\textup{(E.40.3)}, and the literal evaluator
checks $d'=h_2d^\lambda h_2^{-1}$ directly
for each of the $96$ inputs.

For the ternary target use the complete
class-two exponent-three kernel over the
$648$-element quotient as in E.35.
The two successive conjugacy equations retain
the full homogeneous centralizer spaces,
now computed for $d=abc$.
All base conjugators in the $648$-element
group are enumerated.  The exact logarithm
is applied only to the already proved
class-two kernel.
Reconstruction of all nine parameters in
each of the $24$ original planes gives
all $216$ inputs.  The complete calculation
gives \textup{(E.40.4)} and the fifteen
literal conjugators.
Here $|d|=36$ and $|z_4|=9$.
Both multiplier residues, and all their
higher lifts modulo $144$, are one modulo
$36$.  Hence the saved identities give
$d'=h_3d^\lambda h_3^{-1}$ in this factor.

The separate evaluator imports the two
group models and source substitution maps,
but neither conjugacy algorithm.  It
evaluates all nineteen literal conjugators,
reconstructs all $96$ binary and $216$
ternary inputs, and checks the full
multiplier equality at every point.
The opposite-prime corrections have
identity images in the respective target
factors by E.24's explicit power maps.
Finally the actual joint target is the
full direct product proved in E.24.
Thus $(h_2,h_3)$ is an element of that
target and gives \textup{(E.40.2)} for
all $864$ points.

The exact algorithm and exhaustive control
are \texttt{e\_quarter\_torsion\_root.py},
executed in its three modes
\texttt{control}, \texttt{binary}, and
\texttt{ternary}; the separate evaluator is
\texttt{e\_quarter\_torsion\_words\_root.py}.
Their outputs and complete indexed JSON
certificates are retained.  The finite
outcomes keep every original CM fibre.
\end{proof}
